% !TeX program = lualatex
% HTML (lwarp) variant of FTProblems.tex. Sources of the HTML pipeline live in html/; html/build-html.ps1 copies this file
% to the project root for the build and removes it (and everything lwarp makes) afterwards.
\documentclass{ProblemBookHTML}
\addbibresource{d:/Projects/LaTeX/MyPackages/bib/MyBase.bib}

\title{Класична електродинаміка. Збірник задач}
\author{В.~І.~Жданов, С.~М.~Пономаренко, В.~Б.~Долгошей}
\begin{document}
\maketitle
\tableofcontents
% the preface is a page of its own (prepare_html.py turns its \chapter* into \schapterpage)
\ForceHTMLPage
\htmlinput{intro}

\newcommand{\ChaptersOne}{EMstatics,Radiation,SpecialRelativity}
\newcommand{\ChaptersTwo}{EMstaticsInMedia,QuasiStationarMedia,WavesInLinearMedia}

\setcounter{chapter}{-1}
\htmlinput{Vanaliz/Vanaliz}

\part{Мікроскопічна теорія}
\epigraph{\Annabelle Неможливо позбутися відчуття, що ці математичні формули існують незалежно від нас і володіють власним розумом, що вони мудріші за нас, мудріше навіть тих, хто їх відкрив, і що ми дістаємо з них більше, ніж спочатку було закладено\ldots}{Heinrich Hertz}
\htmlinput{Additions/MaxwellEquation}
\multiinclude{\ChaptersOne}[]

\part{Електродинаміка суцільних середовищ}
\epigraph{\Annabelle Вважаю, що більш приземлені та матеріальні науки аж ніяк не можуть бути зневажені у порівнянні з піднесеним вивченням розуму і духу\ldots}{James Clerk Maxwell}
\htmlinput{Additions/MaxwellEquationMacro}
\multiinclude{\ChaptersTwo}[]

% answers to the problems (written by the chapters, as in FTProblems.tex)
\answers
\multiinclude{\ChaptersOne}[-Answers]
\multiinclude{\ChaptersTwo}[-Answers]

% appendices (as in FTProblems.tex)
\appendix
\htmlinput{Additions/Vanaliz}
\htmlinput{Additions/PhysConstants}

% bibliography: the two lists of FTProblems.tex under one heading
\nocite{LL2,ZhdanovFT,ZhdanovESS,ZhdanovRelativity,Fedorchenko1992,Suhakov1975,Jackson1998,GriffithsElectro,GreinerElectrodynamics,LandauLifshitz1975,LandauLifshitz1984,Grechko,BatyginToptygin1978en}
\ForceHTMLPage
\schapter{Література}
\defbibheading{subheading}[\bibname]{\section*{#1}}
\addtocategory{Textbooks}{LL2,ZhdanovFT,ZhdanovESS,ZhdanovRelativity,Fedorchenko1992,Suhakov1975,Jackson1998,GriffithsElectro,GreinerElectrodynamics,LandauLifshitz1975,LandauLifshitz1984}
\addtocategory{Problems}{Grechko,BatyginToptygin1978en}
\printbibliography[category=Textbooks,heading=subheading,title={Підручники та посібники}]
\printbibliography[category=Problems,heading=subheading,title={Задачники}]
\end{document}
